% Standalone resolution manuscript generated from solution.md.
% Compile with XeLaTeX; author and review metadata are maintained in Markdown.
\documentclass[10pt,a4paper]{article}
\usepackage[margin=24mm,headheight=15pt]{geometry}
\usepackage{fontspec}
\setmainfont{DejaVuSerif.ttf}[BoldFont=DejaVuSerif-Bold.ttf,ItalicFont=DejaVuSerif-Italic.ttf,BoldItalicFont=DejaVuSerif-BoldItalic.ttf]
\setsansfont{DejaVuSans.ttf}[BoldFont=DejaVuSans-Bold.ttf,ItalicFont=DejaVuSans-Oblique.ttf]
\setmonofont{DejaVuSansMono.ttf}
\usepackage{amsmath,amssymb,unicode-math}
\setmathfont{latinmodern-math.otf}
\usepackage{xcolor,microtype,parskip,needspace,fancyhdr}
\usepackage{bookmark}
\usepackage{xurl}
\hypersetup{colorlinks=true,urlcolor=blue,linkcolor=blue,pdftitle={RA-10:
constant-loss nuclear-error transfer without matrix
ordering},pdfauthor={George Stepaniants},pdflang={en-GB}}
\urlstyle{same}
\setlength{\emergencystretch}{3em}
\providecommand{\tightlist}{\setlength{\itemsep}{0pt}\setlength{\parskip}{0pt}}
\providecommand{\pandocbounded}[1]{#1}
\setcounter{secnumdepth}{0}
\allowdisplaybreaks[1]
\widowpenalty=10000
\clubpenalty=10000
\pagestyle{fancy}
\fancyhf{}
\fancyhead[L]{\small\sffamily Verified resolution / RA-10}
\fancyhead[R]{\small\sffamily 11 September 2026}
\fancyfoot[L]{\parbox[b]{0.9\textwidth}{\fontsize{8}{10}\selectfont Independent
Codex-agent review; not external human peer review or formal
certification.}}
\fancyfoot[R]{\footnotesize\thepage}
\begin{document}
{\raggedright\Large\bfseries RA-10: constant-loss nuclear-error transfer
without matrix ordering\par}
\vspace{0.4em}
{\normalsize\bfseries George Stepaniants\par}
{\small Department of Computing and Mathematical Sciences, California
Institute of Technology, Pasadena, California, USA\par}

\vspace{0.6em}
\pagestyle{plain}

\emph{11 September 2026}

\textbf{Verification status: independent-agent PASS.} This manuscript
was developed with substantial ChatGPT/Codex assistance. A separate
Codex agent independently checked the complete mathematical argument and
the exact canonical target.
\href{https://github.com/ajt60gaibb/OpenProblemsInNLA/blob/main/references/stepaniants-ra10-2026-09-11/verification/RA-10-independent-review.md}{Detailed
independent review}. This is independent automated-agent review, not
external human peer review or formal verification. No historical
priority or optimality of the constant is claimed. The
\href{https://github.com/ajt60gaibb/OpenProblemsInNLA/blob/main/references/stepaniants-ra10-2026-09-11/README.md}{submission
record} preserves the frozen reviewed draft and its hashes.

\subsection{1. Statement and notation}\label{statement-and-notation}

For a real symmetric matrix \(M\), write \(M_+\) for its positive part
and \(\|M\|_*\) for its nuclear norm. If \(P\) is an orthogonal
projection and \(B=PBP\) is positive semidefinite, write \(f(B)_P\) for
the matrix obtained by applying \(f\) on \(\operatorname{range}(P)\) and
setting its action on \(\ker(P)\) to zero. In particular,
\(f(B)_P=f(B)\) when \(f(0)=0\). The selected projection \(P\) is
retained even when \(B\) has rank smaller than
\(\operatorname{rank}(P)\).

\textbf{Theorem 1.} The exact statement of
\href{https://github.com/ajt60gaibb/OpenProblemsInNLA/blob/main/randomized-and-low-rank-approximation/RA-10/README.md}{RA-10}
holds with the universal constant \(C=11\). More explicitly, let
\(n\ge2\), \(1\le k<n\), let \(A\) and \(\widehat A\) be real symmetric
positive semidefinite matrices, and choose any ordered orthonormal
eigenbasis of \(\widehat A\). Let \(P\) project onto its first \(k\)
vectors and let \(B=\widehat A_k\). For every continuous
operator-monotone \(f:[0,\infty)\to[0,\infty)\) and \(\varepsilon\ge0\),

\[
\|A-B\|_*\le(1+\varepsilon)\|A-A_k\|_*
\]

implies

\[
\|f(A)-f(\widehat A)_k\|_*
\le(1+11\varepsilon)\|f(A)-f(A)_k\|_*.
\]

Both truncations of \(\widehat A\) use the same selected eigenvectors.
No matrix ordering, commutation, invertibility, spectral simplicity, or
strict gap at \(k\) is assumed.

Put \(a_1\ge\cdots\ge a_n\ge0\) for the eigenvalues of \(A\) and
\(\tau=\sum_{i>k}a_i\). The quantity

\[
e(B)=\|A-B\|_*-\tau
\]

is nonnegative by the best-rank-\(k\) approximation property of the
nuclear norm. We first prove the estimate for the ridge atoms
\(f_s(x)=x/(s+x)\), \(s>0\). The proof only uses finite-dimensional
matrix algebra, standard eigenvalue comparison, and the positive
integral representation of operator-monotone functions.

\newpage

\subsection{2. A compression estimate for ridge
functions}\label{a-compression-estimate-for-ridge-functions}

\textbf{Lemma 2.} Let \(A\succeq0\), let \(P\) be an orthogonal
projection of rank \(k\), and put \(C=PAP\). Let \(H\) be the
restriction of \(C\) to \(\operatorname{range}(P)\). For arbitrary
\(s,c>0\), set

\[
D=C-A,\qquad
 d=\operatorname{tr}\bigl((cI_k-H)_+\bigr),\qquad
 g=\frac1{s+c}.
\]

Then

\[
\operatorname{tr}\bigl((f_s(C)-f_s(A))_+\bigr)
\le2g\bigl[\operatorname{tr}(D_+)+d\bigr].\tag{1}
\]

This lemma does not require the eigenvalues of \(C\) to match any
eigenvalues of \(A\).

\textbf{Proof.} First suppose \(H\succ0\). In the orthogonal
decomposition \(\operatorname{range}(P)\oplus\ker(P)\), write

\[
A=\begin{bmatrix}H&E\\E^T&F\end{bmatrix},
\qquad C=\begin{bmatrix}H&0\\0&0\end{bmatrix}.
\]

Positivity of \(A\) gives \(F\succeq E^TH^{-1}E\). Put
\(Z=f_s(C)-f_s(A)\). Let \(v=(u,w)\) be a unit eigenvector of \(Z\) with
positive eigenvalue \(\lambda\). Since \(0\preceq f_s(A)\) and
\(f_s(C)\prec I\), we have \(0<\lambda<1\). Define

\[
\mu=\frac{\lambda}{1+\lambda},\qquad 0<\mu<\frac12.
\]

The resolvent identity \(Z=s[(sI+A)^{-1}-(sI+C)^{-1}]\) and
\(Zv=\lambda v\) give

\[
(sI+A)^{-1}v=(sI+C)^{-1}v+\frac\lambda s v.
\]

Multiplying by \(sI+A\) and using the two block rows yields

\[
Ew=-\mu(sI+H)u,\tag{2}
\]

\[
E^T(sI+\mu H)(sI+H)^{-1}u+(F+\mu sI)w=0.\tag{3}
\]

Taking the inner product of (3) with \(w\) and substituting (2) gives

\[
w^TFw=\mu s(\|u\|^2-\|w\|^2)+\mu^2u^THu.\tag{4}
\]

On the other hand, the Schur-complement inequality and (2) give

\[
\begin{aligned}
w^TFw
&\ge\mu^2u^T(sI+H)H^{-1}(sI+H)u\\
&=\mu^2u^THu+2\mu^2s\|u\|^2
  +\mu^2s^2u^TH^{-1}u.
\end{aligned}
\]

Comparing with (4), and dividing by \(\mu s>0\), proves

\[
\|w\|^2\le(1-2\mu)\|u\|^2-\mu s\,u^TH^{-1}u.
\]

Because \(\|u\|^2+\|w\|^2=1\), it follows that

\[
\|u\|^2\ge\frac1{2(1-\mu)}.\tag{5}
\]

A direct calculation from (2) and (4) gives

\[
v^TDv=\mu s+\mu(2-\mu)u^THu.\tag{6}
\]

Let \(\Delta=(cI_k-H)_+\), so \(H\succeq cI_k-\Delta\). By (5)-(6),

\[
\begin{aligned}
v^TDv
&\ge\mu s+\frac{\mu(2-\mu)c}{2(1-\mu)}
       -\mu(2-\mu)u^T\Delta u\\
&\ge\frac{\lambda(s+c)}2-\mu(2-\mu)u^T\Delta u.
\end{aligned}\tag{7}
\]

\Needspace{6\baselineskip}

The second inequality follows because the difference between its two
scalar terms is

\[
\frac{\mu\,[s(1-2\mu)+c(1-\mu)]}{2(1-\mu)}\ge0.
\]

Take an orthonormal basis of the positive eigenspace of \(Z\) and sum
(7). For its orthogonal projection \(Q\), the variational formula
\(\operatorname{tr}(D_+)=\max_{0\preceq R\preceq I}\operatorname{tr}(RD)\)
gives \(\operatorname{tr}(QD)\le\operatorname{tr}(D_+)\). Also
\(0<\mu(2-\mu)<1\) for each positive eigenvalue, and the
\(P\)-components \(u\) of these orthonormal vectors satisfy
\(\sum uu^T\preceq I_k\). Therefore

\[
\frac{s+c}{2}\operatorname{tr}(Z_+)
\le\operatorname{tr}(D_+)+\operatorname{tr}(\Delta).
\]

This is (1).

If \(H\) is singular, apply the proved case to \(A_r=A+rP\) and
\(C_r=C+rP\), \(r>0\). Their difference remains \(D\), and
\(A_r\succeq0\). For fixed \(s>0\), all resolvents and positive-part
traces are continuous as \(r\) decreases to zero, and
\(\operatorname{tr}((cI-H-rI)_+)\) tends to \(d\). Taking that limit
proves (1) without an invertibility assumption. \(\square\)

\subsection{3. Matching the leading
spectrum}\label{matching-the-leading-spectrum}

Assume \(\tau>0\), so \(c=a_k>0\), and continue to write \(g=1/(s+c)\).
Let \(B_0\) be any positive semidefinite matrix supported on a
rank-\(k\) projection \(P\), whose eigenvalues on that subspace are
\(a_1,\ldots,a_k\). Its eigenvectors are arbitrary. Define

\[
\begin{gathered}
e_0=\|A-B_0\|_*-\tau,\qquad C=PAP,\qquad
H=C\big|_{\operatorname{range}(P)},\\
L=\sum_{i\le k}a_i-\operatorname{tr}H,\qquad
R=\|B_0-C\|_*.
\end{gathered}
\]

The compression eigenvalues \(h_1\ge\cdots\ge h_k\) satisfy
\(0\le h_i\le a_i\) by the min-max principle. Thus \(L\ge0\). Pinching
\(A-B_0\) into the \(P\) and \(I-P\) blocks cannot increase its nuclear
norm, and the complementary \(A\) block is positive semidefinite.
Consequently

\[
e_0\ge R+L.\tag{8}
\]

Put \(e_C=\|A-C\|_*-\tau\). The triangle inequality gives

\[
e_C\le e_0+R.\tag{9}
\]

The trace identity
\(\|M\|_*=2\operatorname{tr}(M_+)-\operatorname{tr}M\) for symmetric
\(M\), applied to \(C-A\), gives

\[
e_C=L+2\operatorname{tr}((C-A)_+).\tag{10}
\]

Moreover,

\[
d=\operatorname{tr}((cI_k-H)_+)
 =\sum_i(c-h_i)_+\le\sum_i(a_i-h_i)=L.\tag{11}
\]

Let \(\tau_s=\sum_{i>k}f_s(a_i)\), and define

\[
L_s=\sum_{i\le k}f_s(a_i)-\operatorname{tr}f_s(H).
\]

\Needspace{8\baselineskip}

For \(a\ge c\) and \(b\ge0\), the exact scalar identity

\[
|f_s(a)-f_s(b)|
=\frac{s|a-b|}{(s+a)(s+b)}\le g|a-b|\tag{12}
\]

holds. Since \(h_i\le a_i\), it implies \(0\le L_s\le gL\).

By Lemma 2, (10)-(11), and the trace-positive-part identity,

\[
\begin{aligned}
\|f_s(A)-f_s(C)\|_*-\tau_s
&=L_s+2\operatorname{tr}((f_s(C)-f_s(A))_+)\\
&\le gL+4g\bigl[\operatorname{tr}((C-A)_+)+d\bigr]\\
&\le g(2e_C+3L).
\end{aligned}\tag{13}
\]

On \(\operatorname{range}(P)\), \(B_0\succeq cI_k\) and \(H\succeq0\).
The resolvent identity and the ideal property of the nuclear norm imply

\[
\begin{aligned}
\|f_s(B_0)-f_s(C)\|_*
&=\bigl\|s(sI+B_0)^{-1}(B_0-C)(sI+C)^{-1}\bigr\|_*\\
&\le gR.
\end{aligned}\tag{14}
\]

Here the calculation is performed on \(\operatorname{range}(P)\); both
matrix-function differences vanish on its orthogonal complement. The
factors have operator norms at most \(1/(s+c)\) and \(1/s\),
respectively. No commutation of \(B_0\) and \(C\) is used.

Combining (8)-(9) and (13)-(14),

\[
\begin{aligned}
\|f_s(A)-f_s(B_0)\|_*-\tau_s
&\le g(2e_C+3L+R)\\
&\le g(2e_0+3R+3L)\\
&\le5g e_0.
\end{aligned}\tag{15}
\]

\subsection{4. Arbitrary selected
approximants}\label{arbitrary-selected-approximants}

Return to \(B=\widehat A_k\) and its prescribed projection \(P\). Let
\(b_1\ge\cdots\ge b_k\ge0\) be its selected eigenvalues. In the same
selected eigenvectors define \(B_0\) with eigenvalues
\(a_1,\ldots,a_k\), and put

\[
r=\sum_{i\le k}|a_i-b_i|=\|B_0-B\|_*.
\]

The eigenvalue perturbation inequality for Hermitian matrices in the
nuclear norm gives

\[
r+\tau\le\|A-B\|_*,\qquad\text{hence }r\le e(B).\tag{16}
\]

This is the standard inequality
\(\sum_i|\lambda_i(X)-\lambda_i(Y)|\le\|X-Y\|_*\) for ordered Hermitian
eigenvalues; the remaining \(n-k\) eigenvalues of \(B\) are zero. For
completeness, put \(J=X-Y\), \(J_-=(-J)_+\), and \(U=X+J_-=Y+J_+\). Then
\(U\succeq X,Y\). The min-max principle gives \(u_i\ge\max(x_i,y_i)\)
for their ordered eigenvalues, so

\[
\begin{aligned}
\sum_i|x_i-y_i|
&\le2\operatorname{tr}U-\operatorname{tr}X-\operatorname{tr}Y\\
&=\operatorname{tr}J_++\operatorname{tr}J_-=\|X-Y\|_*.
\end{aligned}
\]

This proves the inequality, including all repeated-eigenvalue cases.
Also,

\[
e_0\le e(B)+r.\tag{17}
\]

Because \(B_0\) and \(B\) use the same eigenvectors, (12) gives

\[
\|f_s(B_0)-f_s(B)\|_*\le gr.\tag{18}
\]

\Needspace{7\baselineskip}

Equations (15)-(18) prove the absolute excess bound

\[
\begin{aligned}
\|f_s(A)-f_s(B)\|_*-\tau_s
&\le5g\,[e(B)+r]+gr\\
&\le11g\,e(B).
\end{aligned}\tag{19}
\]

Each tail eigenvalue is at most \(c\), so

\[
\tau_s=\sum_{i>k}\frac{a_i}{s+a_i}
\ge\frac\tau{s+c}=g\tau.\tag{20}
\]

Thus, whenever \(e(B)\le\varepsilon\tau\),

\[
\|f_s(A)-f_s(B)\|_*\le(1+11\varepsilon)\tau_s.\tag{21}
\]

In particular, (21) is uniform in \(s>0\) and uses the actual selected
eigenspaces.

\subsection{5. Positive integral combinations and boundary
cases}\label{positive-integral-combinations-and-boundary-cases}

We use the standard positive integral representation of continuous
nonnegative operator-monotone functions on \([0,\infty)\):

\[
f(x)=\alpha+\beta x+\int_{(0,\infty)}\frac{x}{s+x}\,d\nu(s),\tag{22}
\]

where \(\alpha=f(0)\ge0\), \(\beta\ge0\), \(\nu\) is a positive measure
and \(\int_{(0,\infty)}(1+s)^{-1}\,d\nu(s)<\infty\). One primary
reference is P. Chansangiam, \emph{Integral Representations and
Decompositions of Operator Monotone Functions on the Nonnegative Reals},
\href{https://arxiv.org/pdf/1304.7936v1}{arXiv:1304.7936v1}, Proposition
1.1, p.~2. It represents \(f(x)\) as the integral of \(x(1+s)/(x+s)\)
against a finite positive measure \(m\) on the compactified half-line
\([0,\infty]\). Its endpoint atoms give \(\alpha=m(\{0\})\) and
\(\beta=m(\{\infty\})\); on \((0,\infty)\), set
\(d\nu(s)=(1+s)\,dm(s)\), which gives (22) and its stated integrability.
Continuity at zero gives the stated extension there. The original target
requires monotonicity for all real symmetric sizes; realification of
complex Hermitian matrices gives the equivalent hypothesis for the usual
representation theorem.

For fixed finite-dimensional \(A\) and \(B\), (22) may be applied
spectrally. The resulting matrix integrals converge absolutely: each
positive matrix integrand has nuclear norm bounded by the finite sum of
its scalar spectral values, whose integrals are finite. By the triangle
inequality for the integral,

\[
\begin{aligned}
\|f(A)-f(B)_P\|_*
&\le\alpha(n-k)+\beta\|A-B\|_*\\
&\quad+\int_{(0,\infty)}\|f_s(A)-f_s(B)\|_*\,d\nu(s).
\end{aligned}
\]

The best-rank-\(k\) tail of \(f(A)\) is exactly

\[
\tau_f=\alpha(n-k)+\beta\tau
       +\int_{(0,\infty)}\tau_s\,d\nu(s).
\]

If \(\tau>0\), apply (19)-(20), or equivalently (21), to every ridge
atom. The linear term satisfies the original bound with constant one;
the constant term contributes exactly \(\alpha(n-k)\). Since
\(\varepsilon\ge0\),

\[
\|f(A)-f(B)_P\|_*\le(1+11\varepsilon)\tau_f.
\]

\Needspace{12\baselineskip}

Finally, if \(\tau=0\), the hypothesis forces \(\|A-B\|_*=0\), so
\(A=B\) and \(\operatorname{range}(A)\subseteq\operatorname{range}(P)\).
The selected functional-calculus truncation is \(f(B)_P=f(A)P\), and

\[
f(A)-f(B)_P=f(0)(I-P).
\]

Its nuclear norm is \((n-k)f(0)\), exactly the best-rank-\(k\) tail of
\(f(A)\). This handles all zero-tail, zero-matrix and rank-deficient
cases without division by a vanishing quantity. As
\(f(B)_P=f(\widehat A)_k\) by the fixed selected eigenvectors, Theorem 1
follows. \(\square\)

\begingroup
\sloppy

\subsection{References and scope}\label{references-and-scope}

\begin{enumerate}
\def\labelenumi{\arabic{enumi}.}
\tightlist
\item
  D. Persson, R. A. Meyer and C. Musco, \emph{Algorithm-agnostic
  low-rank approximation of operator monotone matrix functions}, SIAM
  Journal on Matrix Analysis and Applications \textbf{46} (2025), 1-21;
  \href{https://arxiv.org/html/2311.14023v2}{arXiv:2311.14023v2},
  Section 1.2 and the closing paragraph of Section 5. Their ordered
  result and counterexamples to constant one are credited; the present
  target is the unordered nuclear-norm question.
\item
  P. Chansangiam, \emph{Integral Representations and Decompositions of
  Operator Monotone Functions on the Nonnegative Reals},
  \href{https://arxiv.org/pdf/1304.7936v1}{arXiv:1304.7936v1} (30 April
  2013), Proposition 1.1, p.~2; Theorem 1.2 gives the equivalent
  weighted-harmonic-mean formulation.
\item
  The
  \href{https://github.com/ajt60gaibb/OpenProblemsInNLA/blob/main/randomized-and-low-rank-approximation/RA-10/README.md}{canonical
  RA-10 entry} and Matthew J. Colbrook's independently reviewed
  \href{https://github.com/ajt60gaibb/OpenProblemsInNLA/blob/main/references/colbrook-transfer-2026-09-11/manuscripts/06_commuting_schatten_transfer.pdf}{commuting
  partial result} are retained in the repository. That result proves the
  necessary lower bound \(C\ge2\). The present proof claims only
  existence with \(C=11\), not sharpness or a Frobenius/Schatten
  generalization.
\end{enumerate}

\endgroup

The substantive new estimate is Lemma 2. The sharper isospectral ridge
inequality and square inequality explored in earlier notes are not
premises of this proof. No numerical search is used to establish the
theorem.
\end{document}
